% Computer Science A --- Java language rules, types, and platform (JVM / JRE / JDK).
% \input from \texttt{main.tex} before \texttt{cs-a/procedural.tex}.

\runningheader{Computer Science A}{Java language \& platform}{Page \thepage\ of \numpages}

\uplevel{\par\textbf{Variables, types, and operators}\par}

\question[4]
What is the difference between assigning and declaring a variable?
\makeemptybox{1.15in}
\begin{solution}
  Declaring introduces the variable (name and type) in the program; assigning stores a value in an already-declared variable (or updates it).
\end{solution}

\question[4]
What is a variable?
\makeemptybox{1in}
\begin{solution}
  A named location in memory that holds a value of a particular type; the value can change during execution.
\end{solution}

\question[4]
What is a primitive datatype? Please list a few and explain them.
\makeemptybox{1.6in}
\begin{solution}
  A type built into Java stored directly as a value (not an object reference). Examples: \texttt{int} (whole numbers), \texttt{double} (floating-point), \texttt{boolean} (true/false), \texttt{char} (single character).
\end{solution}

\question[4]
Which datatype represents text in Java?
\makeemptybox{0.85in}
\begin{solution}
  \texttt{String} (reference type used for text).
\end{solution}

\question[4]
What are shorthand assignment operators?
\makeemptybox{1.1in}
\begin{solution}
  Operators that combine an operation with assignment, e.g.\ \texttt{+=}, \texttt{-=}, \texttt{*=}, \texttt{/=}, \texttt{\%=}, as in \texttt{x += 3} meaning \texttt{x = x + 3}.
\end{solution}

\question[4]
How do you use increment and decrement operators?
\makeemptybox{1.1in}
\begin{solution}
  \texttt{++} adds 1 and \texttt{--} subtracts 1; prefix (\texttt{++x}) vs postfix (\texttt{x++}) differ in the value of the expression when used inside a larger expression.
\end{solution}

\question[4]
What is the modulo operator? How is it useful?
\makeemptybox{1.2in}
\begin{solution}
  \texttt{\%} gives the remainder after division (e.g.\ \texttt{17 \% 5} is \texttt{2}). Useful for wrap-around, parity (even/odd), cycling indices, etc.
\end{solution}

\newpage

\uplevel{\par\textbf{Methods and program structure}\par}

\question[4]
What is a method?
\makeemptybox{1.1in}
\begin{solution}
  A named block of code that can take parameters, perform actions, and optionally return a value; it organizes behavior and can be called (invoked) when needed.
\end{solution}

\question[4]
Explain the main method.
\makeemptybox{1.35in}
\begin{solution}
  \texttt{public static void main(String[] args)} is the entry point the JVM calls when the program starts; it must have that signature (for a standard Java application).
\end{solution}

\question[4]
What is the difference in syntax between calling a method and creating a method?
\makeemptybox{1.25in}
\begin{solution}
  Defining/creating uses a header with return type, name, parameters, and a body in braces; calling uses the method name followed by parentheses with arguments (if any), often as a statement or within an expression.
\end{solution}

\question[4]
Give an example of why you would need to create a method other than the \texttt{main} method.
\makeemptybox{1.2in}
\begin{solution}
  Reuse (same logic in several places), readability (name a subtask), testing one piece of logic, or separating concerns so \texttt{main} stays a thin coordinator.
\end{solution}

\question[4]
Can you describe some of the basic entities of a Java program?
\makeemptybox{1.45in}
\begin{solution}
  Classes, methods, fields/variables, statements, expressions, packages, and (for apps) a class with \texttt{main}; code lives in \texttt{.java} files compiled to bytecode.
\end{solution}

\newpage

\uplevel{\par\textbf{Syntax, blocks, and flow of control}\par}

\question[4]
What do a pair of opening and closing curly braces represent?
\makeemptybox{1in}
\begin{solution}
  A block: a compound statement grouping zero or more statements; used for class bodies, method bodies, and blocks scoped to loops/\texttt{if}/\texttt{switch}.
\end{solution}

\question[4]
What is a conditional statement?
\makeemptybox{1.05in}
\begin{solution}
  A statement that chooses which code runs based on a boolean condition, e.g.\ \texttt{if}, \texttt{if}/\texttt{else}, or \texttt{switch}.
\end{solution}

\question[4]
Explain the different kinds of loops you can create and use in a program.
\makeemptybox{1.5in}
\begin{solution}
  Common Java loops: \texttt{while} (test before each iteration), \texttt{do}-\texttt{while} (test after), \texttt{for} (compact init/test/update), and enhanced \texttt{for} over arrays/collections.
\end{solution}

\question[4]
What is a flow control statement?
\makeemptybox{1.1in}
\begin{solution}
  Statements that alter sequential execution: conditionals, loops, and jumps such as \texttt{break}, \texttt{continue}, and \texttt{return}.
\end{solution}

\question[4]
When would you use an \texttt{if} statement over a \texttt{switch} statement?
\makeemptybox{1.2in}
\begin{solution}
  When conditions are ranges, compound booleans, or non-discrete values; when checking unrelated conditions; or when \texttt{switch} does not fit the value type / pattern (pre--Java~21 style).
\end{solution}

\question[4]
When would you use a \texttt{switch} statement over an \texttt{if} statement?
\makeemptybox{1.2in}
\begin{solution}
  When comparing one expression against many constant discrete values (e.g.\ same integer or \texttt{String} cases) and a multi-branch structure reads clearer than a long \texttt{if}-\texttt{else} chain.
\end{solution}

\newpage

\uplevel{\par\textbf{JDK, JVM, compilation, and the Java platform}\par}

\question[4]
What is the difference between the JDK, JRE, and JVM?
\makeemptybox{1.45in}
\begin{solution}
  JVM executes bytecode; JRE (historically) bundled JVM + libraries for running apps; JDK includes the compiler (\texttt{javac}), tools, and libraries for developing (and usually running) Java programs.
\end{solution}

\question[4]
What is the difference between source code and bytecode?
\makeemptybox{1.15in}
\begin{solution}
  Source code is human-readable \texttt{.java} text; bytecode is the platform-independent intermediate form (\texttt{.class}) the JVM executes after compilation.
\end{solution}

\question[4]
How would you describe the compilation process for Java?
\makeemptybox{1.35in}
\begin{solution}
  The compiler translates \texttt{.java} source into \texttt{.class} bytecode; the JVM loads classes and interprets/JIT-compiles bytecode to machine code at runtime.
\end{solution}

\question[4]
What are some of the benefits of Java?
\makeemptybox{1.35in}
\begin{solution}
  Portability (``write once, run anywhere'' on a JVM), strong ecosystem, garbage collection, mature tooling, wide use in industry/education, and rich standard libraries.
\end{solution}

\newpage

\uplevel{\par\textbf{Arrays}\par}

\question[4]
What is an array? Why is it useful?
\makeemptybox{1.25in}
\begin{solution}
  A fixed-size ordered collection indexed from zero holding elements of one type; useful for storing many related values without separate variables per item.
\end{solution}

\question[4]
How can you iterate over an array?
\makeemptybox{1.15in}
\begin{solution}
  Indexed \texttt{for} loop with an index from \texttt{0} to \texttt{length - 1}, or an enhanced \texttt{for} (\texttt{for (T x : arr)}).
\end{solution}

\question[4]
How would you access the last value in an array if you do not know the size of the array?
\makeemptybox{1in}
\begin{solution}
  Index with \texttt{array.length - 1}. If the array may be empty, check \texttt{length > 0} first.
\end{solution}

\newpage

\uplevel{\par\textbf{Input, \texttt{Scanner}, and strings}\par}

\question[4]
What is the \texttt{Scanner} class used for? Give an example of how you would use it.
\makeemptybox{1.45in}
\begin{solution}
  Parsing primitive types and strings from an input stream (e.g.\ keyboard). Example: \texttt{Scanner s = new Scanner(System.in); int n = s.nextInt();}
\end{solution}

\question[4]
If you received text input from the user, how would you go about comparing it to a value, like \texttt{"yes"} or \texttt{"no"}?
\makeemptybox{1.15in}
\begin{solution}
  Compare with \texttt{String} methods \texttt{equals} or \texttt{equalsIgnoreCase}; do not use \texttt{==} on reference types for content equality.
\end{solution}

\newpage

\uplevel{\par\textbf{Debugging}\par}

\question[4]
What is a stack trace? How can you use it to debug your application?
\makeemptybox{1.35in}
\begin{solution}
  A report of active stack frames at an exception (or thread dump), showing call order and line numbers; use it to find which method threw and the chain of calls that led there.
\end{solution}

\question[4]
What steps would you take to debug your program if it crashed with an error message in the console?
\makeemptybox{1.45in}
\begin{solution}
  Read the exception type and message, find the topmost frame in your code, inspect that line, reproduce reliably, then fix logic or add guards/logging.
\end{solution}

\question[4]
What steps would you take to debug your program if it runs, but gives you the wrong results?
\makeemptybox{1.45in}
\begin{solution}
  Trace expected vs actual with small inputs, use logging/breakpoints, narrow which step first diverges, verify assumptions and formulas, check types and edge cases.
\end{solution}

\question[4]
If there was an error in the console when you tried to run your program, how would you debug?
\makeemptybox{1.35in}
\begin{solution}
  Same general approach as a crash: read the error, locate the failing line from the stack trace, fix the cause (syntax, null, bounds, logic), recompile, and retest.
\end{solution}

\newpage

\uplevel{\par\textbf{Packages, scope, static vs.\ instance}\par}

\question[4]
What is a package and why would we use one?
\makeemptybox{1.2in}
\begin{solution}
  A namespace for related classes (folder-like hierarchy); avoids name clashes, organizes large projects, and controls visibility with access modifiers.
\end{solution}

\question[4]
What is a fully qualified class name?
\makeemptybox{1.05in}
\begin{solution}
  The package plus simple class name, e.g.\ \texttt{java.util.Scanner}, uniquely identifying the class on the classpath.
\end{solution}

\question[4]
What are the different scopes in Java?
\makeemptybox{1.35in}
\begin{solution}
  Common scopes: class/static fields, instance fields, local variables in a method or block, and parameters; inner classes/lambdas add nested lexical scope rules.
\end{solution}

\question[4]
If I define a variable within a method, how can I access its value outside of the method?
\makeemptybox{1.05in}
\begin{solution}
  You cannot access a local variable outside its declaring block; return it, store it in a field, or pass it via parameters to code that needs it.
\end{solution}

\question[4]
What is the difference between using an instance variable and a static variable?
\makeemptybox{1.35in}
\begin{solution}
  Each object has its own instance fields; static fields belong to the class and are shared by all instances (one storage location per class loader).
\end{solution}

\question[4]
What is the difference between calling an instance method and a static method?
\makeemptybox{1.35in}
\begin{solution}
  Instance methods run on a specific object and can use \texttt{this} and instance state; static methods belong to the class, are called via the class name, and cannot access instance fields/\texttt{this} directly.
\end{solution}

\newpage

\uplevel{\par\textbf{Computing vocabulary}\par}

\question[4]
What is an operating system?
\makeemptybox{1.15in}
\begin{solution}
  System software that manages hardware, processes, memory, filesystems, and devices, and provides services so applications can run (e.g.\ scheduling, I/O).
\end{solution}

\question[4]
What does the term ``full stack'' mean?
\makeemptybox{1.15in}
\begin{solution}
  Typically end-to-end development across layers: front end (UI), back end (server, APIs, database), and deployment---``full stack developer'' works across those areas.
\end{solution}

\newpage

\begingradingrange{csalang} 

\uplevel{\par\textbf{Java platform (JVM, JRE, JDK)}\par}

\question[4]
Java uses a virtual machine (JVM) to execute code, providing platform independence.
\begin{choices}
  \CorrectChoice TRUE
  \choice FALSE
\end{choices}
\begin{solution}
  Bytecode runs on the JVM, which is implemented for each platform so the same \texttt{.class} files can run cross-platform.
\end{solution}

\question[4]
The code that the developer writes is also known as \blank[5em].
\begin{choices}
  \CorrectChoice Source code
  \choice Binary
  \choice Bytecode
  \choice An application
\end{choices}
\begin{solution}
  Human-readable Java source is compiled to bytecode for the JVM.
\end{solution}

\question[4]
In order to run Java code, you need the \blank[3em].
\begin{choices}
  \choice JVM
  \CorrectChoice JRE
  \choice JIT
  \choice JDK
\end{choices}
\begin{solution}
  \textbf{Deployment view:} end users need a \textbf{JRE} (JVM plus standard libraries such as \texttt{java.base}) so bytecode can run against the Java API. A bare JVM without those libraries cannot run a typical application.

  \textbf{Roles:}
  \begin{itemize}
    \item \textbf{JVM:} executes bytecode.
    \item \textbf{JRE:} runnable environment (JVM + standard library modules).
    \item \textbf{JDK:} developer kit (compiler \texttt{javac}, tools, plus a JRE for development).
  \end{itemize}

  \textbf{Distribution:} \texttt{jlink} builds a custom runtime image with only needed modules; \texttt{jpackage} bundles an app with a runtime into installers.

  \begin{center}
  \small
  \begin{tabular}{@{}p{0.42\linewidth}p{0.48\linewidth}@{}}
  \hline
  \textbf{If you want to\ldots} & \textbf{You need\ldots} \\
  \hline
  Run a Java program & JRE (includes the JVM) \\
  Develop a Java program & JDK (includes the JRE and JVM) \\
  Understand the core engine & JVM \\
  \hline
  \end{tabular}
  \end{center}
\end{solution}

\question[4]
How does Java compile your code?
\begin{choices}
  \choice Source code converted to bytecode by JRE; then the bytecode is converted to machine code by the JDK
  \choice Source code is converted to bytecode by the JDK; then the bytecode is compiled line by line via the JRE; then bytecode is converted to machine code
  \choice Source code is converted to bytecode by the JRE; then the bytecode is converted line by line using the JDK's JIT; the JIT converts the bytecode to machine code
  \CorrectChoice Source code is converted to bytecode by the JDK; then JIT compilation is performed by the JVM which processes the bytecode line by line for execution
\end{choices}
\begin{solution}
  \texttt{javac} (JDK) produces bytecode; the JVM loads and executes it, with the JIT compiling frequently executed code to native instructions.
\end{solution}

\newpage

\uplevel{\par\textbf{Language syntax, types, and program structure}\par}

\question[4]
What is the purpose of the \texttt{static} keyword in Java?
\begin{choices}
  \choice It makes a variable constant
  \CorrectChoice It indicates that a method belongs to the class, not to an instance of the class
  \choice To make a variable mutable
  \choice It prevents a class from being extended
\end{choices}
\begin{solution}
  \texttt{static} fields and methods are shared by the class; call static members with \texttt{ClassName.member} without an instance.
\end{solution}

\question[4]
The \blank[5em]\ keyword is used for sharing a variable or a method that is the same for every instance; it belongs to the class rather than to an instance of the class.
\begin{choices}
  \choice \lstinline!final!
  \choice \lstinline!new!
  \CorrectChoice \lstinline!static!
  \choice \lstinline!abstract!
\end{choices}
\begin{solution}
  \texttt{static} members are per-class; \texttt{final} controls reassignment/immutability patterns; \texttt{new} creates instances; \texttt{abstract} marks incomplete members or types.
\end{solution}

\question[4]
Which of the following is the correct way to declare a variable in Java?
\begin{choices}
  \choice \lstinline!variable name;!
  \choice \lstinline!declare name;!
  \CorrectChoice \lstinline!int name;!
  \choice None of the above
\end{choices}
\begin{solution}
  \texttt{variable} and \texttt{declare} are not types; \lstinline!int name;! declares (but does not initialize) an \texttt{int}.
\end{solution}

\question[4]
How should you compare two strings for equality in Java?
\begin{choices}
  \choice Using the \lstinline!==! operator
  \CorrectChoice Using the \lstinline!equals()! method
  \choice Using the \lstinline!compare()! method
  \choice Using the \lstinline!compareTo()! method
\end{choices}
\begin{solution}
  \texttt{==} compares references; \texttt{equals} compares character content. \texttt{compare}/\texttt{compareTo} relate to ordering, not value equality.
\end{solution}

\question[4]
How do you separate statements in Java?
\begin{choices}
  \choice \lstinline!,!
  \choice \lstinline!.!
  \CorrectChoice \lstinline!;!
  \choice None of the above
\end{choices}
\begin{solution}
  Statements end with a semicolon; blocks use braces without semicolons after the closing brace (except in special cases).
\end{solution}

\question[4]
An expression can be made up of variables, constants, and/or operators.
\begin{choices}
  \CorrectChoice TRUE
  \choice FALSE
\end{choices}
\begin{solution}
  Expressions evaluate to values: literals, variables, operators, and method calls combine into larger expressions.
\end{solution}

\question[4]
Is \lstinline!com._12example! a valid package name?
\begin{choices}
  \CorrectChoice TRUE
  \choice FALSE
\end{choices}
\begin{solution}
  Segments may contain underscores; a segment must not \emph{start} with a digit---here \texttt{\_12example} starts with underscore, which is allowed.
\end{solution}

\newpage

\question[4]
There is no need to create an object for the class to access a static variable or static method.
\begin{choices}
  \CorrectChoice TRUE
  \choice FALSE
\end{choices}
\begin{solution}
  Access via \texttt{ClassName.staticMember}; an instance is not required (though \texttt{instance.static} is legal but discouraged).
\end{solution}

\question[4]
How do you get the size of an array \lstinline!myArray!?
\begin{choices}
  \choice \lstinline!myArray.size!
  \CorrectChoice \lstinline!myArray.length!
  \choice \lstinline!myArray.size()!
  \choice \lstinline!myArray.length()!
\end{choices}
\begin{solution}
  Arrays use the \texttt{length} field; \texttt{List} and similar use \texttt{size()}.
\end{solution}

\question[4]
In which scope is \texttt{x} declared?
\begin{lstlisting}[style=javaexam]
public class MyClass {
  // What is the variable scope for x?
  static String x = "test string 1";
  String y = "test string 2";
  void perform() {
    String z = "test string 4";
    for (int i = 0; i < 4; i++) {
      System.out.println(i);
    }
  }
}
\end{lstlisting}
\begin{choices}
  \CorrectChoice Static (class) scope
  \choice Instance scope
  \choice Method scope
  \choice Block scope
\end{choices}
\begin{solution}
  \texttt{static} fields belong to the class, not to any one instance; \texttt{y} is instance scope, \texttt{z} is method-local, \texttt{i} is loop block scope.
\end{solution}

\question[4]
Which of the following is \emph{not} a primitive type in Java?
\begin{choices}
  \CorrectChoice \lstinline!String!
  \choice \lstinline!int!
  \choice \lstinline!double!
  \choice \lstinline!short!
\end{choices}
\begin{solution}
  \texttt{String} is a reference type (class); the eight primitives include \texttt{int}, \texttt{double}, \texttt{short}, etc.
\end{solution}

\question[4]
What are the benefits of packages?
\begin{choices}
  \choice Packages enable you to hide your classes
  \choice Packages help make code visible to all other parts of the program
  \CorrectChoice Allow you to group like classes and interfaces together to assist with organizing your projects
  \choice Allows developers to write their code in a file
\end{choices}
\begin{solution}
  Packages provide namespaces, organize source trees, and help avoid name collisions across libraries.
\end{solution}

\question[4]
What is the correct syntax for the \texttt{main} method in Java?
\begin{choices}
  \choice \lstinline!public static main(String[] args) {}!
  \choice \lstinline!public static int main(String[] args) {}!
  \choice \lstinline!static void main(String[] args) {}!
  \CorrectChoice \lstinline!public static void main(String[] args) {}!
\end{choices}
\begin{solution}
  The entry point must be \texttt{public static void main(String[] args)} (or \texttt{String... args}).
\end{solution}

\newpage

\uplevel{\par\textbf{JVM memory: heap and stack}\par}
% Threads, reflection, and the Stream API are covered in \texttt{cs-a/distributed\_computing.tex}.

\question[4]
New objects are always created in \blank[8em].
\begin{choices}
  \choice Stack memory
  \choice The \texttt{main} method
  \CorrectChoice Heap memory
  \choice The string pool only
\end{choices}
\begin{solution}
  Object instances (and arrays) live on the \textbf{heap}; local variables holding references reside in stack frames. The string pool holds interned literals/interned strings, not all objects.
\end{solution}

\question[4]
The heap is where method invocations and their local frames are stored.
\begin{choices}
  \choice TRUE
  \CorrectChoice FALSE
\end{choices}
\begin{solution}
  Each thread's \textbf{call stack} holds frames for active method invocations (locals, operand stack, etc.). The \textbf{heap} stores object data; the JVM moves between bytecode instructions and the heap through references in those frames.
\end{solution}

\newpage

\uplevel{\par\textbf{Types, wrappers, and memory (practice bank)}\par}

\question[4]
You can assign a variable of a larger capacity to a smaller variable without casting.
\begin{choices}
  \choice TRUE
  \CorrectChoice FALSE
\end{choices}
\begin{solution}
  Narrowing primitive assignments require an explicit cast in Java (e.g.\ \texttt{int} to \texttt{byte}).
\end{solution}

\question[4]
Which of the following is \emph{not} a wrapper class?
\begin{choices}
  \choice \texttt{Boolean}
  \choice \texttt{Long}
  \choice \texttt{Double}
  \choice \texttt{Short}
  \CorrectChoice All of these are wrapper classes
\end{choices}
\begin{solution}
  Each listed type has a wrapper in \texttt{java.lang}; the intended ``none of the above'' answer is that they are \emph{all} wrappers.
\end{solution}

\question[4]
What does \emph{not} belong in the Java heap (in the usual CS-A teaching model)?
\begin{choices}
  \CorrectChoice Methods
  \choice Objects
  \choice String pool
  \choice None of the above
\end{choices}
\begin{solution}
  Method bytecode lives in \textbf{Metaspace}/method area (not the object heap). Object instances and the string pool (for interned strings) are heap-related concepts in coursework.
\end{solution}

\endgradingrange{csalang}
